\chapter{サイコ・サイバネティクス}
\epigraph{%
  \itshape ``The guidance unit does not project the target's future;
  it continuously computes the angular tracking error,
  generating corrective fin deflections until divergence reaches zero.''
}{%
  \textbf{AIM-9 Sidewinder Technical Guidance \& Control Manual}\\
  Naval Weapons Center, China Lake
}
\section{サイバネティクスからサイコ・サイバネティクスへ：目標追従系としての確信$\Omega^{\mathrm{ego}}_{\mathrm{predict:profit}}$と目標回避系としての確信$\Omega^{\mathrm{ego}}_{\mathrm{predict:threat}}$}

\subsection{サイバネティクスの歴史的系譜：目的論的フィードバック制御}
サイバネティクス（Cybernetics）の歴史は、第二次世界大戦の最中、MITの数学者ノーバート・ウィーナー（Norbert Wiener）らによって着手された防空火器管制の数理解析に遡る。
高速で予測困難な軌道を描く敵爆撃機を撃墜するためには、現在の航空機の位置だけでなく、射撃手や制御機構の「予測」と「偏差のフィードバック修正」を統合した数理モデルが不可欠であった。

1948年、ウィーナーは名著『サイバネティクス：動物と機械における制御と通信』を上梓し、機械工学、生理学、通信理論を横断する画期的な概念を提唱した。
その核心は、\textbf{「系（システム）が外部からのセンサー情報を取り込み、目標値との誤差（エラー）を検出し、ネガティブ・フィードバック（負帰還）をかけることによって自律的に目標へと収束していく閉ループの目的論的メカニズム」}である。

同時期に開催されたメイシー会議（1946–1953年）には、ウィーナーをはじめ、神経生理学者ウォーレン・マカロック、人類学者グレゴリー・ベイトソン、数学者ジョン・フォン・ノイマンらが集い、生体と機械に共通する自律制御の原理が議論された。
ここで確立された「サーボメカニズム（自動制御系）」の概念は、単なる兵器や自動機械の域を超え、生命現象や認知システムの動作原理を解き明かす普遍的パラダイムへと急速に拡大していった。

\subsection{モルツの「サイコ・サイバネティクス」と自己イメージ}
この工学的サーボメカニズムの概念を、人間の精神・臨床心理学へと応用したのが、形成外科医マクスウェル・モルツ（Maxwell Maltz）である。
1960年、モルツは『サイコ・サイバネティクス（Psycho-Cybernetics）』を発表した。

美容整形手術を手がけていたモルツは、顔面の醜い傷や奇形を完全に修復したにもかかわらず、ある患者たちは依然として「自分は醜く、誰からも愛されない」と信じ込み、手術前とまったく同じ劣等感や引きこもり行動を続けるという奇怪な現象に直面した。
外科用メスで肉体をいくら物理的に変容させても、心の中の振る舞いが1ミリも変わらない人間が存在する一方で、わずかな処置で劇的に人生を好転させる患者もいた。

この臨床観察から、モルツはひとつの結論を導き出した。
人間の脳・神経系には、工学的な自動操縦装置（オートパイロット）とまったく同一のサーボメカニズムが内蔵されている。
そして、そのサーボ機構が追い続ける\textbf{「設定目標値（セットポイント）」こそが、個人の内面に強固に刻み込まれた「自己イメージ（Self-Image）」}である、と。
脳はこの設定目標値を絶対的な基準とし、現実の行動や感情との間にズレが生じると、ネガティブ・フィードバックを激しく作動させて、何が何でも自己イメージ通りの状態へと生体を引き戻そうとするのである。

\subsection{サイコ・サイバネティクスシステムの実体：目標値としての確信 $\Omega^{\mathrm{ego}}$}

本稿が構築してきた神経系疑似ベイズ推論機のフレームワークにおいて、モルツが直観した「サイコ・サイバネティクスシステムの設定目標値」の正体は、まさに\textbf{profitとなるマクロなEgoの未来予測の確信 $\Omega^{\mathrm{ego}}_{\mathrm{predict:profit}}$} そのものとして物理的に定式化される。

\begin{equation}
    \text{Target Set-point} \equiv \Omega^{\mathrm{ego}}_{\mathrm{predict:profit}}
\end{equation}

目標追従系（Attractor System）としてのサイコ・サイバネティクス機構とは、感覚ノイズや内的想起 $e_{\mathrm{noise}}$ を受容した際、事前予期 $\hat{\Omega}^{\mathrm{ego}}$ に従って確定したナラティブ解 $\Omega^{\mathrm{ego}}$ を「不変の参照基準（誘因目標値）」としてレジスタに固定し、全身の身体運動 $\lambda$、情動反応 $\Lambda$、および内部対話 $A_d$ をその目標値へ向けて閉ループ制御するサーボ機構にほかならない。

\begin{equation}
    \text{Error}_{\mathrm{track}} = \| \text{Current State} - \Omega^{\mathrm{ego}}_{\mathrm{predict:profit}} \| \xrightarrow{\text{Feedback Control}} \Delta (\lambda, \Lambda, A_d) \longrightarrow 0
\end{equation}

一方、Egoの未来予期の確信が「脅威（threat）」として確定した場合（$\Omega^{\mathrm{ego}}_{\mathrm{predict:threat}}$）、系は目標追従型から\textbf{「目標回避系（Repulsive Control System）」}へと切り替わる。
この場合、制御器は脅威状態との距離を安全マージン $D_{\mathrm{safe}}$ 以上に引き離す、あるいは脅威ポテンシャル場からの反発力（回避誤差）をゼロに抑え込むように作用する。

\begin{align}
    \text{Error}_{\mathrm{avoid}} &= \max\left( 0,\, D_{\mathrm{safe}} - \| \text{Current State} - \Omega^{\mathrm{ego}}_{\mathrm{predict:threat}} \| \right) \\
                                  &\xrightarrow{\text{Feedback Control}} \Delta (\lambda, \Lambda, A_d) \longrightarrow 0
\end{align}

この目標回避制御のフィードバック出力こそが、生体レベルにおける「逃走（Flight）または闘争（Fight）」の防衛行動として具現化する。

\subsection{破滅的ナラティブへの収束と自己破壊的ホメオスタシス}

ここで極めて重要なパラドックスが立ち現れる。
\textbf{確定するナラティブの確信 $\Omega^{\mathrm{ego}}$ は、必ずしも生体に快楽や幸福をもたらす目標値（profit）として設定されるわけではない。むしろ多くの人間において、それは「エゴの不安や破滅の予測」そのものを追従目標値として固定してしまう。}

なぜなら、生物の進化史において神経系に課された至上命題は「幸福になること」ではなく「死を避けること（生存確率の最大化と不確実性の極小化）」だからである。
生命体にとって最も破滅的な状態とは、外部環境の脅威が未確定なまま放置され、推論ループが際限なく空転し続けること（$\boldsymbol{e}_{\mathrm{seq}} \xrightarrow[]{\hat{\Omega}^{\mathrm{ego}}} \mid$）である。何が起きるかわからない宙ぶらりんの不確定状態は、莫大な計算資源を浪費し、即座の防衛行動を麻痺させる。

したがって、DMNを主軸とするEgoの計算モジュールは、\textbf{「たとえ自身にとって悲惨で苦痛に満ちた結末であっても、不確定な未来を放置するよりは、最悪のストーリーを確信 $\Omega^{\mathrm{ego}}$ として確定させてしまった方が、系の自由エネルギー（不確実性）を最小化できる」}と判断する。

\begin{itemize}
    \item 「私はどうせ誰からも愛されない」という確信 $\Omega^{\mathrm{ego}}_{\mathrm{narrative:rejection}}$
    \item 「努力しても必ず最後には裏切られる」という確信 $\Omega^{\mathrm{ego}}_{\mathrm{narrative:failure}}$
    \item 「世界は悪意に満ちており、私は被害者である」という確信 $\Omega^{\mathrm{ego}}_{\mathrm{narrative:victim}}$
\end{itemize}

これらの破壊的ナラティブは、心理的には耐え難い苦痛をもたらす。
しかし、サイバネティクス制御の観点から見れば、ひとたびこの確信 $\Omega^{\mathrm{ego}}_{\mathrm{narrative}}$ が設定目標値（アトラクター）としてレジスタに固定されさえすれば、推論系としては極めて「安定」した状態を獲得する。
「愛されない」「失敗する」という結末が確定しているならば、もはや未知の事態に怯える必要はなく、ただその設定目標値と現実のズレを解消するように行動（人間関係の自己破壊的サボタージュや過剰な防衛攻撃 $\Lambda$）をフィードバック制御すればよいからである。

病的な自己否定や慢性的な不安に苦しむ人間が、なぜそこから抜け出せないのか。
それは彼らが快楽を求めているからではない。彼らの神経系疑似ベイズ推論機が、エゴの破滅的ナラティブを「絶対的な確信 $\Omega^{\mathrm{ego}}$」としてセットポイントに据え、その破滅という目標を寸分の狂いもなく忠実に達成し続ける「完全無欠なサイコ・サイバネティクスサーボ系」として、あまりにも正常に機能してしまっているからなのである。


